% ŚCIĄGA — przykłady użycia szablonu. Usuń % ŚCIĄGA — przykłady użycia szablonu. Usuń % ŚCIĄGA — przykłady użycia szablonu. Usuń % ŚCIĄGA — przykłady użycia szablonu. Usuń \include{rozdzialy/99-przyklady} z main.tex przed oddaniem.
\chapter{Przykłady użycia szablonu (do usunięcia)}
\label{chap:przyklady}

\section{Tekst, cytowania i odwołania}

Zwykły akapit jest wyjustowany, ma wcięcie 1,27~cm i~interlinię 1,5. Cytowanie pojedynczej
pozycji wygląda tak~\cite{przyklad-ksiazka}, kilku naraz tak~\cite{przyklad-artykul,
przyklad-konferencja,przyklad-www}. Kolejność w~bibliografii wynika z~kolejności pierwszego
cytowania. Odwołania do innych miejsc: rozdział~\ref{chap:wstep}, rysunek~\ref{fig:przyklad},
tabela~\ref{tab:przyklad}, wzór~\eqref{eq:przyklad}, listing~\ref{lst:java}\footnote{Przypis
dolny wstawia się poleceniem \texttt{\textbackslash footnote}.}.

Jednoliterowe spójniki i~przyimki (a, i, o, u, w, z) wiąże się z~następnym słowem tyldą:
\texttt{w\textasciitilde{}pracy}, aby nie zostawały na końcu wiersza. Liczby z~przecinkiem
dziesiętnym: \num{1234.5} albo z~jednostką: \SI{12.5}{\milli\second}.

\section{Listy}

\begin{itemize}
  \item pierwszy element listy punktowanej,
  \item drugi element:
  \begin{itemize}
    \item element drugiego poziomu,
    \item kolejny element drugiego poziomu,
  \end{itemize}
  \item trzeci element.
\end{itemize}

\begin{enumerate}
  \item pierwszy krok,
  \item drugi krok.
\end{enumerate}

\section{Rysunek, wykres i tabela}

Rysunek wstawiany jest z~katalogu \texttt{grafika/} (najlepiej w~PDF lub PNG). Podpis jest pod
rysunkiem, numeracja jest ciągła w~całej pracy.

\begin{figure}[htbp]
  \centering
  \includegraphics[width=0.5\textwidth]{example-image}
  \caption{Przykładowy rysunek}
  \label{fig:przyklad}
\end{figure}

Wykresy mają osobną numerację (w~oryginale „Wykres 1”, „Wykres 2”, …). Jeśli wolisz traktować
je jak rysunki, użyj zwykłego \texttt{figure}.

\begin{wykres}[htbp]
  \centering
  \includegraphics[width=0.5\textwidth]{example-image-a}
  \caption{Przykładowy wykres}
  \label{wyk:przyklad}
\end{wykres}

Tabela w~stylu oryginału: pełna siatka, szary wiersz nagłówka, wyśrodkowane komórki, podpis pod
tabelą (aby przenieść podpis nad tabelę, przesuń \texttt{\textbackslash caption} przed
\texttt{tabular}).

\begin{table}[htbp]
  \centering
  \begin{tabular}{|c|c|c|}
    \hline
    \rowcolor{tabnaglowek}
    \textbf{Kolumna 1} & \textbf{Kolumna 2} & \textbf{Kolumna 3} \\
    \hline
    A & \num{0.95} & \num{120} \\
    \hline
    B & \num{0.87} & \num{340} \\
    \hline
  \end{tabular}
  \caption{Przykładowa tabela}
  \label{tab:przyklad}
\end{table}

\section{Wzory i listingi}

\begin{equation}
  F_1 = 2 \cdot \frac{\mathit{precision} \cdot \mathit{recall}}{\mathit{precision} + \mathit{recall}}
  \label{eq:przyklad}
\end{equation}

\begin{lstlisting}[language=Java, caption={Przykładowy kod w~języku Java}, label={lst:java}]
public class Main {
    public static void main(String[] args) {
        // komentarz z polskimi znakami: zażółć gęślą jaźń
        System.out.println("Witaj!");
    }
}
\end{lstlisting}

\begin{lstlisting}[language=Rust, caption={Przykładowy kod w~języku Rust}, label={lst:rust}]
fn main() {
    let liczby: Vec<u64> = (1..=10).collect();
    let suma: u64 = liczby.iter().sum();
    println!("Suma: {}", suma);
}
\end{lstlisting}

Kod można też wczytać wprost z~pliku źródłowego:

\begin{verbatim}
\lstinputlisting[language=Rust, caption={Opis}]{../rust/src/main.rs}
\end{verbatim}
 z main.tex przed oddaniem.
\chapter{Przykłady użycia szablonu (do usunięcia)}
\label{chap:przyklady}

\section{Tekst, cytowania i odwołania}

Zwykły akapit jest wyjustowany, ma wcięcie 1,27~cm i~interlinię 1,5. Cytowanie pojedynczej
pozycji wygląda tak~\cite{przyklad-ksiazka}, kilku naraz tak~\cite{przyklad-artykul,
przyklad-konferencja,przyklad-www}. Kolejność w~bibliografii wynika z~kolejności pierwszego
cytowania. Odwołania do innych miejsc: rozdział~\ref{chap:wstep}, rysunek~\ref{fig:przyklad},
tabela~\ref{tab:przyklad}, wzór~\eqref{eq:przyklad}, listing~\ref{lst:java}\footnote{Przypis
dolny wstawia się poleceniem \texttt{\textbackslash footnote}.}.

Jednoliterowe spójniki i~przyimki (a, i, o, u, w, z) wiąże się z~następnym słowem tyldą:
\texttt{w\textasciitilde{}pracy}, aby nie zostawały na końcu wiersza. Liczby z~przecinkiem
dziesiętnym: \num{1234.5} albo z~jednostką: \SI{12.5}{\milli\second}.

\section{Listy}

\begin{itemize}
  \item pierwszy element listy punktowanej,
  \item drugi element:
  \begin{itemize}
    \item element drugiego poziomu,
    \item kolejny element drugiego poziomu,
  \end{itemize}
  \item trzeci element.
\end{itemize}

\begin{enumerate}
  \item pierwszy krok,
  \item drugi krok.
\end{enumerate}

\section{Rysunek, wykres i tabela}

Rysunek wstawiany jest z~katalogu \texttt{grafika/} (najlepiej w~PDF lub PNG). Podpis jest pod
rysunkiem, numeracja jest ciągła w~całej pracy.

\begin{figure}[htbp]
  \centering
  \includegraphics[width=0.5\textwidth]{example-image}
  \caption{Przykładowy rysunek}
  \label{fig:przyklad}
\end{figure}

Wykresy mają osobną numerację (w~oryginale „Wykres 1”, „Wykres 2”, …). Jeśli wolisz traktować
je jak rysunki, użyj zwykłego \texttt{figure}.

\begin{wykres}[htbp]
  \centering
  \includegraphics[width=0.5\textwidth]{example-image-a}
  \caption{Przykładowy wykres}
  \label{wyk:przyklad}
\end{wykres}

Tabela w~stylu oryginału: pełna siatka, szary wiersz nagłówka, wyśrodkowane komórki, podpis pod
tabelą (aby przenieść podpis nad tabelę, przesuń \texttt{\textbackslash caption} przed
\texttt{tabular}).

\begin{table}[htbp]
  \centering
  \begin{tabular}{|c|c|c|}
    \hline
    \rowcolor{tabnaglowek}
    \textbf{Kolumna 1} & \textbf{Kolumna 2} & \textbf{Kolumna 3} \\
    \hline
    A & \num{0.95} & \num{120} \\
    \hline
    B & \num{0.87} & \num{340} \\
    \hline
  \end{tabular}
  \caption{Przykładowa tabela}
  \label{tab:przyklad}
\end{table}

\section{Wzory i listingi}

\begin{equation}
  F_1 = 2 \cdot \frac{\mathit{precision} \cdot \mathit{recall}}{\mathit{precision} + \mathit{recall}}
  \label{eq:przyklad}
\end{equation}

\begin{lstlisting}[language=Java, caption={Przykładowy kod w~języku Java}, label={lst:java}]
public class Main {
    public static void main(String[] args) {
        // komentarz z polskimi znakami: zażółć gęślą jaźń
        System.out.println("Witaj!");
    }
}
\end{lstlisting}

\begin{lstlisting}[language=Rust, caption={Przykładowy kod w~języku Rust}, label={lst:rust}]
fn main() {
    let liczby: Vec<u64> = (1..=10).collect();
    let suma: u64 = liczby.iter().sum();
    println!("Suma: {}", suma);
}
\end{lstlisting}

Kod można też wczytać wprost z~pliku źródłowego:

\begin{verbatim}
\lstinputlisting[language=Rust, caption={Opis}]{../rust/src/main.rs}
\end{verbatim}
 z main.tex przed oddaniem.
\chapter{Przykłady użycia szablonu (do usunięcia)}
\label{chap:przyklady}

\section{Tekst, cytowania i odwołania}

Zwykły akapit jest wyjustowany, ma wcięcie 1,27~cm i~interlinię 1,5. Cytowanie pojedynczej
pozycji wygląda tak~\cite{przyklad-ksiazka}, kilku naraz tak~\cite{przyklad-artykul,
przyklad-konferencja,przyklad-www}. Kolejność w~bibliografii wynika z~kolejności pierwszego
cytowania. Odwołania do innych miejsc: rozdział~\ref{chap:wstep}, rysunek~\ref{fig:przyklad},
tabela~\ref{tab:przyklad}, wzór~\eqref{eq:przyklad}, listing~\ref{lst:java}\footnote{Przypis
dolny wstawia się poleceniem \texttt{\textbackslash footnote}.}.

Jednoliterowe spójniki i~przyimki (a, i, o, u, w, z) wiąże się z~następnym słowem tyldą:
\texttt{w\textasciitilde{}pracy}, aby nie zostawały na końcu wiersza. Liczby z~przecinkiem
dziesiętnym: \num{1234.5} albo z~jednostką: \SI{12.5}{\milli\second}.

\section{Listy}

\begin{itemize}
  \item pierwszy element listy punktowanej,
  \item drugi element:
  \begin{itemize}
    \item element drugiego poziomu,
    \item kolejny element drugiego poziomu,
  \end{itemize}
  \item trzeci element.
\end{itemize}

\begin{enumerate}
  \item pierwszy krok,
  \item drugi krok.
\end{enumerate}

\section{Rysunek, wykres i tabela}

Rysunek wstawiany jest z~katalogu \texttt{grafika/} (najlepiej w~PDF lub PNG). Podpis jest pod
rysunkiem, numeracja jest ciągła w~całej pracy.

\begin{figure}[htbp]
  \centering
  \includegraphics[width=0.5\textwidth]{example-image}
  \caption{Przykładowy rysunek}
  \label{fig:przyklad}
\end{figure}

Wykresy mają osobną numerację (w~oryginale „Wykres 1”, „Wykres 2”, …). Jeśli wolisz traktować
je jak rysunki, użyj zwykłego \texttt{figure}.

\begin{wykres}[htbp]
  \centering
  \includegraphics[width=0.5\textwidth]{example-image-a}
  \caption{Przykładowy wykres}
  \label{wyk:przyklad}
\end{wykres}

Tabela w~stylu oryginału: pełna siatka, szary wiersz nagłówka, wyśrodkowane komórki, podpis pod
tabelą (aby przenieść podpis nad tabelę, przesuń \texttt{\textbackslash caption} przed
\texttt{tabular}).

\begin{table}[htbp]
  \centering
  \begin{tabular}{|c|c|c|}
    \hline
    \rowcolor{tabnaglowek}
    \textbf{Kolumna 1} & \textbf{Kolumna 2} & \textbf{Kolumna 3} \\
    \hline
    A & \num{0.95} & \num{120} \\
    \hline
    B & \num{0.87} & \num{340} \\
    \hline
  \end{tabular}
  \caption{Przykładowa tabela}
  \label{tab:przyklad}
\end{table}

\section{Wzory i listingi}

\begin{equation}
  F_1 = 2 \cdot \frac{\mathit{precision} \cdot \mathit{recall}}{\mathit{precision} + \mathit{recall}}
  \label{eq:przyklad}
\end{equation}

\begin{lstlisting}[language=Java, caption={Przykładowy kod w~języku Java}, label={lst:java}]
public class Main {
    public static void main(String[] args) {
        // komentarz z polskimi znakami: zażółć gęślą jaźń
        System.out.println("Witaj!");
    }
}
\end{lstlisting}

\begin{lstlisting}[language=Rust, caption={Przykładowy kod w~języku Rust}, label={lst:rust}]
fn main() {
    let liczby: Vec<u64> = (1..=10).collect();
    let suma: u64 = liczby.iter().sum();
    println!("Suma: {}", suma);
}
\end{lstlisting}

Kod można też wczytać wprost z~pliku źródłowego:

\begin{verbatim}
\lstinputlisting[language=Rust, caption={Opis}]{../rust/src/main.rs}
\end{verbatim}
 z main.tex przed oddaniem.
\chapter{Przykłady użycia szablonu (do usunięcia)}
\label{chap:przyklady}

\section{Tekst, cytowania i odwołania}

Zwykły akapit jest wyjustowany, ma wcięcie 1,27~cm i~interlinię 1,5. Cytowanie pojedynczej
pozycji wygląda tak~\cite{przyklad-ksiazka}, kilku naraz tak~\cite{przyklad-artykul,
przyklad-konferencja,przyklad-www}. Kolejność w~bibliografii wynika z~kolejności pierwszego
cytowania. Odwołania do innych miejsc: rozdział~\ref{chap:wstep}, rysunek~\ref{fig:przyklad},
tabela~\ref{tab:przyklad}, wzór~\eqref{eq:przyklad}, listing~\ref{lst:java}\footnote{Przypis
dolny wstawia się poleceniem \texttt{\textbackslash footnote}.}.

Jednoliterowe spójniki i~przyimki (a, i, o, u, w, z) wiąże się z~następnym słowem tyldą:
\texttt{w\textasciitilde{}pracy}, aby nie zostawały na końcu wiersza. Liczby z~przecinkiem
dziesiętnym: \num{1234.5} albo z~jednostką: \SI{12.5}{\milli\second}.

\section{Listy}

\begin{itemize}
  \item pierwszy element listy punktowanej,
  \item drugi element:
  \begin{itemize}
    \item element drugiego poziomu,
    \item kolejny element drugiego poziomu,
  \end{itemize}
  \item trzeci element.
\end{itemize}

\begin{enumerate}
  \item pierwszy krok,
  \item drugi krok.
\end{enumerate}

\section{Rysunek, wykres i tabela}

Rysunek wstawiany jest z~katalogu \texttt{grafika/} (najlepiej w~PDF lub PNG). Podpis jest pod
rysunkiem, numeracja jest ciągła w~całej pracy.

\begin{figure}[htbp]
  \centering
  \includegraphics[width=0.5\textwidth]{example-image}
  \caption{Przykładowy rysunek}
  \label{fig:przyklad}
\end{figure}

Wykresy mają osobną numerację (w~oryginale „Wykres 1”, „Wykres 2”, …). Jeśli wolisz traktować
je jak rysunki, użyj zwykłego \texttt{figure}.

\begin{wykres}[htbp]
  \centering
  \includegraphics[width=0.5\textwidth]{example-image-a}
  \caption{Przykładowy wykres}
  \label{wyk:przyklad}
\end{wykres}

Tabela w~stylu oryginału: pełna siatka, szary wiersz nagłówka, wyśrodkowane komórki, podpis pod
tabelą (aby przenieść podpis nad tabelę, przesuń \texttt{\textbackslash caption} przed
\texttt{tabular}).

\begin{table}[htbp]
  \centering
  \begin{tabular}{|c|c|c|}
    \hline
    \rowcolor{tabnaglowek}
    \textbf{Kolumna 1} & \textbf{Kolumna 2} & \textbf{Kolumna 3} \\
    \hline
    A & \num{0.95} & \num{120} \\
    \hline
    B & \num{0.87} & \num{340} \\
    \hline
  \end{tabular}
  \caption{Przykładowa tabela}
  \label{tab:przyklad}
\end{table}

\section{Wzory i listingi}

\begin{equation}
  F_1 = 2 \cdot \frac{\mathit{precision} \cdot \mathit{recall}}{\mathit{precision} + \mathit{recall}}
  \label{eq:przyklad}
\end{equation}

\begin{lstlisting}[language=Java, caption={Przykładowy kod w~języku Java}, label={lst:java}]
public class Main {
    public static void main(String[] args) {
        // komentarz z polskimi znakami: zażółć gęślą jaźń
        System.out.println("Witaj!");
    }
}
\end{lstlisting}

\begin{lstlisting}[language=Rust, caption={Przykładowy kod w~języku Rust}, label={lst:rust}]
fn main() {
    let liczby: Vec<u64> = (1..=10).collect();
    let suma: u64 = liczby.iter().sum();
    println!("Suma: {}", suma);
}
\end{lstlisting}

Kod można też wczytać wprost z~pliku źródłowego:

\begin{verbatim}
\lstinputlisting[language=Rust, caption={Opis}]{../rust/src/main.rs}
\end{verbatim}
